\chapter{A Neuron Is Not One Device}
% full version: chapters/18_eetypes.tex

In the first chapter we sorted neurons by the substance they release. An engineer would sort them differently: by what device each one works as. And it turns out that behind the same appearance hides a whole box of electronic components.

\section*{A Box of Components}

The leaky bucket from the second chapter is a leaky integrator: it adds up signals over a short window. With a very small window it becomes a coincidence detector. A neuron that responds only to changes and falls silent at once is a filter that passes changes. A neuron that runs by itself is a metronome, an oscillator; if an input adjusts its tempo, it is a voltage-controlled oscillator. A neuron that puts out bursts of clicks is a burst generator, writing ``dashes.'' A neuron that clicks in time with a sound wave is a phase detector.

\pencilfig{18}{\pencilart{18}}{Behind the same appearance hides a whole box of components: a resistor, a capacitor, a coil, an oscillator.}

There are also devices we have not talked about before. A \emph{resonator} is like a swing: it responds best to pushes of a certain rhythm, like a swing that goes higher only if you push it at the right moment. Inside the neuron the role of the spring is played by a slow current that resists change. An \emph{integrator without leak} holds a value for a long time, like a capacitor that does not leak; this is how the network that holds the position of the eyes works when you look to the side. But for this the feedback must be tuned almost perfectly, and that is a costly adjustment. A \emph{latch} is a neuron that a short push switches into lasting activity until it is switched off, like a latching switch; in the neurons that control muscles, serotonin turns this mode on.

\section*{A Chameleon Device}

The main discovery of this sorting: these are not different parts but different \emph{modes} of one and the same part. The same neuron can be an integrator by day and a burst generator by night, depending on what the radio stations tell it. An engineer would call the brain a field-programmable gate array: one circuit, while the settings change what the parts do.

That neurons can work in all these modes is measured. How the brain uses mode switching for computation is mostly a hypothesis.

\fullversion{18}
